\chapter{Réalisation de la plateforme clinique et déploiement}
\label{chap:realisation-plateforme}

Les modèles décrits au chapitre précédent ne valent que par l'usage qu'un praticien
peut en faire. Ce chapitre décrit la plateforme qui les met à disposition :
l'interface applicative et ses points d'entrée, les mécanismes de sécurité, la
stratégie d'archivage, l'interface utilisateur et la production du compte rendu. Il se
termine par la conteneurisation, la chaîne d'intégration continue et la campagne de
validation fonctionnelle.

% =============================================================================
\section{Interface applicative}
% =============================================================================

\subsection{Organisation des modules}

L'interface applicative est bâtie sur FastAPI, choisi pour son modèle asynchrone, sa
validation déclarative des données et la génération automatique de sa documentation.
Le code est réparti en modules à responsabilité unique, décrits au
tableau~\ref{tab:modules}.

\begin{table}[H]
\centering
\caption{Modules de l'interface applicative}
\label{tab:modules}
\small
\begin{tabular}{@{}lrL{8.3cm}@{}}
\toprule
\textbf{Module} & \textbf{Lignes} & \textbf{Responsabilité} \\
\midrule
\texttt{main.py}          & 695 & Point d'entrée, routes métier et d'administration,
                                  migration au démarrage, journalisation \\
\texttt{ml\_routes.py}    & 984 & Routes d'inférence : analyse, extraction de
                                  marqueurs, prédiction de sévérité, lecture de
                                  fichiers externes \\
\texttt{ml\_models.py}    & 117 & Définition des trois architectures de réseaux \\
\texttt{preprocessing.py} & 199 & Décodage EDF, résolution des dérivations, calcul
                                  des métriques AASM \\
\texttt{osa\_predictor.py}& 113 & Chargement des modèles de sévérité, prédiction et
                                  attribution des contributions \\
\texttt{db\_models.py}    & 113 & Modèle de données et relations \\
\texttt{schemas.py}       & 163 & Schémas de validation des entrées et sorties \\
\texttt{auth.py}          &  51 & Jetons, empreintes de mots de passe, dépendance
                                  d'authentification \\
\texttt{b2\_storage.py}   & 121 & Client de stockage objet, dépôt et signature des
                                  liens \\
\texttt{database.py}      &  19 & Connexion et session \\
\bottomrule
\end{tabular}
\end{table}

Deux mécanismes transversaux méritent d'être signalés.

Le premier est la \textbf{migration au démarrage}. Le schéma ayant évolué au fil des
sprints, l'application exécute à son lancement une série d'ajouts de colonnes
conditionnels, ce qui permet de déployer une nouvelle version sur une base existante
sans intervention manuelle. Cette approche convient à un projet de cette taille ;
au-delà, un outil de migration versionnée serait préférable, et c'est ce que nous
recommandons en perspective.

Le second est la \textbf{signature des liens à la sérialisation}. Les schémas de
sortie déclarent un validateur qui intercepte toute référence de fichier au moment de
composer la réponse et la remplace par un lien signé temporaire. La conséquence est
importante : aucune route ne peut divulguer une référence brute par inadvertance,
puisque la conversion est appliquée par le schéma et non par la route.

\subsection{Points d'entrée}

L'interface expose trente points d'entrée, regroupés par domaine au
tableau~\ref{tab:endpoints}. La documentation interactive est produite automatiquement
au format OpenAPI.

\begin{table}[H]
\centering
\caption{Points d'entrée de l'interface applicative}
\label{tab:endpoints}
\small
\begin{tabular}{@{}lL{6.6cm}cL{2.6cm}@{}}
\toprule
\textbf{Méthode} & \textbf{Chemin} & \textbf{Accès} & \textbf{Domaine} \\
\midrule
POST & \texttt{/token}                        & public   & \multirow{2}{*}{Authentification} \\
GET  & \texttt{/users/me}                     & connecté & \\
\midrule
POST & \texttt{/admin/doctors}                & admin & \multirow{8}{*}{Administration} \\
GET  & \texttt{/admin/doctors}                & admin & \\
PUT  & \texttt{/admin/doctors/\{id\}/lifecycle}      & admin & \\
POST & \texttt{/admin/doctors/\{id\}/reset-password} & admin & \\
POST & \texttt{/admin/hospitals}              & admin & \\
GET  & \texttt{/admin/hospitals}              & admin & \\
GET  & \texttt{/admin/dashboard-stats}        & admin & \\
GET  & \texttt{/admin/audit-logs}             & admin & \\
GET  & \texttt{/admin/audit-logs/export}      & admin & \\
\midrule
POST & \texttt{/patients}                     & médecin & \multirow{4}{*}{Patients} \\
GET  & \texttt{/patients}                     & médecin & \\
GET  & \texttt{/patients/\{id\}}              & médecin & \\
GET  & \texttt{/doctor/stats}                 & médecin & \\
\midrule
POST & \texttt{/patients/\{id\}/psgs}         & médecin & \multirow{8}{*}{Examens} \\
PUT  & \texttt{/psgs/\{id\}}                  & médecin & \\
POST & \texttt{/psgs/\{id\}/upload\_edf}      & médecin & \\
POST & \texttt{/psgs/\{id\}/upload\_hypnogram} & médecin & \\
POST & \texttt{/psgs/\{id\}/upload\_hypnogram\_annotated} & médecin & \\
POST & \texttt{/psgs/\{id\}/upload\_osa\_report} & médecin & \\
POST & \texttt{/psgs/\{id\}/annotations}      & médecin & \\
GET  & \texttt{/psgs/\{id\}/annotations}      & médecin & \\
\midrule
GET  & \texttt{/doctors}                      & médecin & \multirow{5}{*}{Collaboration} \\
POST & \texttt{/conversations}                & médecin & \\
GET  & \texttt{/conversations}                & médecin & \\
GET  & \texttt{/conversations/psg/\{id\}}     & médecin & \\
GET  & \texttt{/conversations/\{id\}/messages} & médecin & \\
POST & \texttt{/conversations/\{id\}/messages} & médecin & \\
\midrule
GET  & \texttt{/health}                       & public & \multirow{7}{*}{Inférence} \\
POST & \texttt{/channels}                     & public & \\
POST & \texttt{/analyze}                      & public & \\
POST & \texttt{/extract\_features}            & public & \\
POST & \texttt{/predict\_osa}                 & public & \\
POST & \texttt{/predict\_osa\_custom}         & public & \\
POST & \texttt{/parse\_features\_file}        & public & \\
\bottomrule
\end{tabular}
\end{table}

Il faut relever ici une \textbf{faiblesse de la version livrée} : les routes
d'inférence ne sont pas protégées par le contrôle d'accès. Ce choix, initialement fait
pour faciliter les essais pendant le sprint de mise au point, n'a pas été rétabli. Un
appelant non authentifié peut donc soumettre un enregistrement à l'analyse. Aucune
donnée patient n'est exposée — ces routes ne lisent ni n'écrivent en base — mais elles
consomment des ressources de calcul. L'ajout de la dépendance d'authentification, qui
tient en une ligne par route, figure parmi les corrections à apporter avant toute mise
en service réelle.

\capturefig{swagger}{0.9}{Documentation OpenAPI générée automatiquement, avec les
trente points d'entrée regroupés par domaine.}{fig:swagger}

% =============================================================================
\section{Sécurité et conformité}
% =============================================================================

\subsection{Mise en œuvre}

Les quatre mécanismes conçus au chapitre~\ref{chap:conception} se traduisent ainsi.

Les \textbf{mots de passe} sont hachés par bcrypt, dont le sel aléatoire par compte
interdit les tables précalculées et dont le coût algorithmique ralentit une attaque
par force brute. Le condensat seul est stocké ; la vérification recalcule l'empreinte
à partir du mot de passe fourni.

Les \textbf{jetons} sont signés en HMAC-SHA256 et portent l'adresse électronique de
l'utilisateur, son rôle et une date d'expiration fixée à sept jours. La clé de
signature provient de l'environnement d'exécution. Chaque route protégée déclare une
dépendance qui décode le jeton, vérifie la signature et l'échéance, puis charge le
compte correspondant.

Le \textbf{contrôle d'accès} superpose deux niveaux. Le rôle est vérifié en tête de
route ; la propriété de la ressource l'est ensuite, en confrontant l'identifiant du
praticien à celui du patient consulté. Un médecin ne peut donc pas accéder au dossier
d'un confrère, même en devinant un identifiant.

Les \textbf{liens d'accès} aux fichiers sont signés et expirent au bout de vingt-quatre
heures. Le service de stockage vérifie la signature avant de servir l'objet. Une
subtilité a été traitée : les enregistrements et les fichiers tabulaires sont servis
avec une en-tête forçant le téléchargement, tandis que les images et les comptes rendus
restent affichables dans le navigateur.

\subsection{Cycle de vie des licences}

Un cinquième niveau, propre au contexte clinique, conditionne l'ouverture de session.
Trois vérifications s'ajoutent à celle du mot de passe : le compte ne doit pas être
suspendu, il ne doit pas être en attente d'approbation, et sa date d'échéance ne doit
pas être dépassée. Chacune produit un message distinct, de sorte que le praticien sache
s'il doit contacter son administrateur ou renouveler son abonnement.

Ce dispositif permet de \textbf{révoquer un accès sans détruire de données}, ce qui est
la seule manière acceptable de gérer une fin de contrat sur un système qui héberge des
dossiers médicaux.

\subsection{Traçabilité}

Chaque action sensible inscrit dans le journal son auteur, son rôle, son horodatage,
la nature de l'action et l'adresse d'origine. Sont journalisées les connexions, les
créations de compte, les modifications de statut et les réinitialisations de mot de
passe. Le journal est consultable par l'administrateur et exportable au format
tabulaire pour un contrôle externe.

Comme exposé au chapitre~\ref{chap:conception}, la table est délibérément dénormalisée :
elle recopie l'adresse et le rôle plutôt que de référencer le compte, afin que la trace
survive à la suppression de celui-ci.

\capturefig{audit}{0.9}{Journal d'audit : consultation filtrée et export tabulaire.}{fig:audit}

\subsection{Limites au regard du règlement européen}

Le système n'est pas conforme au règlement général sur la protection des données, et il
serait malhonnête de le prétendre. Trois écarts principaux subsistent.

Les \textbf{données patient ne sont pas chiffrées au repos} : elles reposent sur le
chiffrement offert par le service de stockage et par le système de fichiers de la base,
sans couche applicative supplémentaire. Le \textbf{droit à l'effacement} n'est pas
outillé : aucune procédure n'efface un patient et ses fichiers archivés en une
opération. Enfin, aucune \textbf{durée de conservation} n'est définie ni appliquée.

Ces trois points sont identifiés et documentés ; ils relèvent d'un travail de mise en
conformité qui dépasse le périmètre d'un projet de fin d'études, mais qu'une mise en
service réelle imposerait.

% =============================================================================
\section{Archivage sur stockage objet}
% =============================================================================

Les fichiers volumineux sont confiés à un service de stockage objet compatible avec
l'interface S3. Chaque objet est rangé dans un préfixe portant le nom du patient et
nommé selon le motif \texttt{patient\_date\_alea}, ce qui rend l'arborescence lisible
depuis la console du fournisseur sans consulter la base.

Le point remarquable de cette réalisation est le \textbf{transfert en arrière-plan}.
Un enregistrement pèse plusieurs centaines de mégaoctets ; attendre la fin de son
téléversement avant d'afficher les résultats aurait immobilisé le praticien pendant
plusieurs minutes. L'enchaînement retenu est donc le suivant : l'analyse est exécutée,
les résultats sont affichés, l'examen est créé en base, puis le transfert démarre sans
bloquer l'interface. Une vignette flottante affiche la progression de chaque transfert
et disparaît quelques secondes après son achèvement.

Un échec de transfert n'interrompt rien : il est signalé dans la vignette, et l'examen
reste pleinement exploitable puisque les résultats d'analyse sont stockés en base
indépendamment des fichiers. C'est le comportement modélisé par le diagramme d'états du
chapitre~\ref{chap:conception}.

% =============================================================================
\section{Interface utilisateur}
% =============================================================================

\subsection{Architecture des composants}

L'interface est une application monopage React, construite par Vite. Elle est
organisée en composants regroupés par domaine fonctionnel, décrits au
tableau~\ref{tab:composants}.

\begin{table}[H]
\centering
\caption{Principaux composants de l'interface}
\label{tab:composants}
\small
\begin{tabular}{@{}lrL{8.0cm}@{}}
\toprule
\textbf{Composant} & \textbf{Lignes} & \textbf{Rôle} \\
\midrule
\texttt{AnalysisResults} & 1134 & Restitution complète : hypnogramme, métriques,
                                  marqueurs, formulaire clinique, rapport de sévérité \\
\texttt{PatientList}     &  993 & Registre des patients, historique des examens,
                                  documents archivés \\
\texttt{CustomOSA}       &  658 & Prédiction à partir d'une saisie manuelle ou d'un
                                  rapport importé \\
\texttt{Hypnogram}       &  552 & Tracé interactif de l'hypnogramme, de la saturation
                                  et de la densité d'événements \\
\texttt{DoctorList}      &  539 & Administration des comptes et des licences \\
\texttt{DeveloperPipeline} & 497 & Visualisation de la chaîne de traitement \\
\texttt{LandingPage}     &  406 & Page publique et authentification \\
\texttt{AnalysisWizard}  &  428 & Assistant d'analyse en quatre étapes \\
\texttt{HospitalsDashboard} & 233 & Gestion des établissements \\
\texttt{ConsultationsView}  & 198 & Boîte de réception des sollicitations \\
\texttt{AuditLogsDashboard} & 194 & Consultation et export du journal \\
\texttt{CollaborationChat}  & 171 & Fil de discussion rattaché à un examen \\
\bottomrule
\end{tabular}
\end{table}

L'ensemble représente environ \num{7700} lignes. Le choix de composants volumineux
plutôt que d'une décomposition fine se justifie par la cohésion : le composant de
restitution, par exemple, orchestre des visualisations qui partagent toutes le même
état d'analyse, et le fractionner aurait imposé un mécanisme de partage sans bénéfice.

\subsection{Parcours d'analyse}

Le parcours principal se déroule en quatre étapes. Le praticien déclare d'abord le
montage disponible, puis la granularité de classification souhaitée, dépose ensuite le
fichier, et consulte enfin les résultats. Chaque étape n'expose qu'une seule décision,
assortie d'une recommandation par défaut.

\capturefig{wizard-1}{0.86}{Assistant d'analyse, étape 1 : déclaration du montage
disponible.}{fig:wizard1}

\capturefig{wizard-2}{0.86}{Assistant d'analyse, étape 2 : choix de la granularité de
classification.}{fig:wizard2}

Pendant le traitement, un écran de progression affiche les tracés physiologiques et
détaille les cinq étapes du pipeline. Ce retour visuel n'est pas décoratif : une
analyse dure plusieurs dizaines de secondes, et une interface muette pendant ce laps
de temps laisserait croire à un blocage.

\capturefig{progression}{0.86}{Écran de progression : tracés physiologiques et étapes
du pipeline.}{fig:progression}

\subsection{Restitution des résultats}

La page de résultats superpose sur un même axe temporel l'hypnogramme prédit, la courbe
de saturation et la densité d'événements respiratoires. Cette superposition est le
choix de visualisation le plus utile de l'interface : elle permet de constater
directement qu'une désaturation coïncide avec une sortie de sommeil paradoxal, ce que
trois graphiques séparés ne montreraient pas.

\capturefig{resultats}{0.92}{Restitution : hypnogramme, saturation et densité
d'événements sur un axe commun, puis métriques AASM.}{fig:resultats}

Les métriques AASM sont présentées sous forme de cartes, chacune accompagnée de sa
plage de référence et d'une comparaison aux valeurs de la cohorte SHHS. Une valeur hors
norme est signalée visuellement, et une liste d'alertes cliniques reprend en clair les
anomalies détectées.

\subsection{Prédiction de sévérité et explication}

Sous les résultats de classification, un formulaire recueille les données
démographiques et oxymétriques. La prédiction produit une classe de sévérité, ses
probabilités et la décomposition des facteurs, présentés en deux colonnes : ce qui
aggrave l'estimation et ce qui la tempère.

Un second parcours permet une prédiction sans enregistrement, à partir d'une saisie
manuelle ou d'un rapport externe importé au format tabulaire ou balisé. Cette voie
répond à un usage réel : disposer d'un compte rendu produit par un autre laboratoire
sans avoir accès au signal brut.

\capturefig{custom-osa}{0.86}{Prédiction de sévérité à partir d'une saisie manuelle.}{fig:customosa}

\subsection{Multilinguisme et thème sombre}

L'interface est intégralement traduite en français et en anglais. Les chaînes sont
centralisées dans un catalogue unique et la bascule est immédiate, sans rechargement.
Un thème sombre est également disponible, adapté à la lecture nocturne — un praticien
qui interprète un enregistrement le fait rarement en plein jour.

\capturefig{dark}{0.86}{Tableau de bord du médecin en thème sombre.}{fig:dark}

% =============================================================================
\section{Production du compte rendu}
% =============================================================================

Le compte rendu est produit côté client : les sections retenues sont converties en
images par capture du rendu, puis assemblées dans un document paginé. Cette approche
garantit que le document reproduit exactement ce que le praticien a validé à l'écran,
sans qu'un moteur de rendu distinct puisse introduire un écart.

Le praticien choisit les sections à inclure : hypnogramme, métriques AASM, classe de
sévérité, facteurs explicatifs, tracés bruts et comparaison aux normes. Cette
composition à la carte répond à des destinataires différents — un confrère n'attend pas
le même document qu'un dossier d'assurance.

Le document est ensuite archivé au même titre que les autres pièces de l'examen.

\capturefig{rapport}{0.78}{Compte rendu clinique généré, avec les sections
sélectionnées par le praticien.}{fig:rapport}

% =============================================================================
\section{Conteneurisation}
% =============================================================================

La pile est décrite par un fichier d'orchestration unique qui déclare trois services.

Le service de \textbf{base de données} s'appuie sur une image PostgreSQL allégée, avec
un volume persistant et une sonde de disponibilité. Cette sonde est indispensable :
sans elle, l'interface applicative démarrerait avant que la base n'accepte les
connexions et échouerait à créer son schéma.

Le service d'\textbf{interface applicative} est construit à partir d'une image Python
allégée. Son fichier de construction comporte une particularité qui mérite explication :
PyTorch y est installé depuis l'index dédié aux processeurs, \emph{avant} les autres
dépendances. Sans cette précaution, l'installation télécharge plusieurs gigaoctets de
bibliothèques CUDA inutiles en déploiement — l'inférence en production s'exécutant sur
processeur — et fait échouer la construction sur une machine de développement.

Le service de \textbf{serveur web} est construit en deux étapes : une première image
Node compile l'application, une seconde image Nginx ne reçoit que le résultat compilé.
L'image finale ne contient donc ni les sources, ni les dépendances de développement.
Une configuration Nginx spécifique redirige toutes les routes inconnues vers le
document principal, ce qu'exige le routage côté client.

\begin{table}[H]
\centering
\caption{Services de la pile conteneurisée}
\label{tab:services}
\small
\begin{tabular}{@{}lllL{5.4cm}@{}}
\toprule
\textbf{Service} & \textbf{Image de base} & \textbf{Port} & \textbf{Particularité} \\
\midrule
Base de données & \texttt{postgres:15-alpine} & 5432 &
Volume persistant, sonde de disponibilité \\
\addlinespace[2pt]
Interface applicative & \texttt{python:3.10-slim} & 8000 &
PyTorch en version processeur installé en premier \\
\addlinespace[2pt]
Serveur web & \texttt{node:20-alpine} puis \texttt{nginx:alpine} & 80 &
Construction en deux étapes, redirection du routage client \\
\bottomrule
\end{tabular}
\end{table}

L'ensemble démarre par une commande unique, et les secrets — identifiants de base,
clé de signature, accès au stockage — sont injectés par variables d'environnement.

% =============================================================================
\section{Intégration continue}
% =============================================================================

Une chaîne d'intégration continue s'exécute à chaque livraison sur la branche
principale et à chaque proposition de fusion. Elle comporte trois travaux, décrits au
tableau~\ref{tab:ci}.

\begin{table}[H]
\centering
\caption{Travaux de la chaîne d'intégration continue}
\label{tab:ci}
\small
\begin{tabular}{@{}lL{10.4cm}@{}}
\toprule
\textbf{Travail} & \textbf{Étapes} \\
\midrule
Interface applicative &
Installation des dépendances avec PyTorch en version processeur ; contrôle syntaxique
de tous les modules ; démarrage d'une instance PostgreSQL de service ; démarrage réel
de l'application ; vérification du schéma OpenAPI et de la présence des points
d'entrée attendus ; contrôle du cloisonnement des rôles \\
\addlinespace[3pt]
Interface utilisateur &
Installation reproductible des dépendances ; analyse statique ; construction de
production ; publication du bundle en tant qu'artefact \\
\addlinespace[3pt]
Images de conteneur &
Construction des deux images avec cache de couches, conditionnée à la réussite des
deux travaux précédents \\
\bottomrule
\end{tabular}
\end{table}

Le travail portant sur l'interface applicative ne se contente pas d'un contrôle
syntaxique : il \textbf{démarre réellement l'application} contre une base de données
éphémère, puis exerce trois vérifications de comportement. Un appel sans jeton doit
recevoir un refus d'authentification ; une connexion administrateur doit délivrer un
jeton ; et ce jeton d'administrateur, présenté au registre des patients, doit recevoir
un refus d'autorisation. Ces trois assertions couvrent l'essentiel du cloisonnement
des rôles et échoueraient immédiatement si une dépendance de sécurité venait à
disparaître au fil d'une modification.

\capturefig{ci}{0.9}{Exécution de la chaîne d'intégration continue.}{fig:ci}

% =============================================================================
\section{Validation fonctionnelle}
% =============================================================================

En l'absence de tests automatisés exhaustifs, la validation a reposé sur une campagne
de scénarios fonctionnels exécutés manuellement à chaque fin de sprint. Le
tableau~\ref{tab:tests} en présente les principaux.

\begin{table}[H]
\centering
\caption{Scénarios de validation fonctionnelle}
\label{tab:tests}
\small
\begin{tabular}{@{}clL{7.2cm}c@{}}
\toprule
\textbf{N\textsuperscript{o}} & \textbf{Scénario} & \textbf{Critère de réussite} &
\textbf{État} \\
\midrule
1 & Connexion valide & Jeton délivré, espace de travail ouvert & \checkmark \\
2 & Connexion refusée & Compte suspendu ou licence expirée rejeté avec message
    distinct & \checkmark \\
3 & Cloisonnement des rôles & Un médecin n'atteint aucune route d'administration &
    \checkmark \\
4 & Cloisonnement des patients & Un médecin n'accède pas au dossier d'un confrère &
    \checkmark \\
5 & Analyse complète & Enregistrement déposé, hypnogramme et métriques affichés &
    \checkmark \\
6 & Montage incomplet & Dérivation absente : repli appliqué, analyse aboutie &
    \checkmark \\
7 & Fichier invalide & Format non reconnu refusé avec message explicite & \checkmark \\
8 & Prédiction de sévérité & Classe, probabilités et facteurs restitués & \checkmark \\
9 & Import de rapport externe & Formulaire prérempli depuis un fichier tabulaire ou
    balisé & \checkmark \\
10 & Archivage en arrière-plan & Résultats affichés avant la fin du transfert &
    \checkmark \\
11 & Échec d'archivage & Transfert interrompu : examen consultable, erreur signalée &
    \checkmark \\
12 & Lien expiré & Lien réutilisé après échéance : accès refusé & \checkmark \\
13 & Annotation d'époque & Note enregistrée et retrouvée au rechargement & \checkmark \\
14 & Consultation confraternelle & Message émis par un praticien, reçu par l'autre &
    \checkmark \\
15 & Compte rendu & Document produit avec les seules sections sélectionnées &
    \checkmark \\
16 & Journal d'audit & Actions sensibles tracées et exportées & \checkmark \\
17 & Bascule linguistique & Aucune chaîne non traduite sur le parcours principal &
    \checkmark \\
18 & Démarrage de la pile & Trois services opérationnels en une commande & \checkmark \\
\bottomrule
\end{tabular}
\end{table}

Il faut reconnaître la limite de cette approche : une campagne manuelle n'est ni
reproductible ni exhaustive, et elle ne protège pas contre les régressions entre deux
exécutions. Les trois assertions automatisées de la chaîne d'intégration continue
constituent un premier filet de sécurité, mais une véritable couverture de tests —
unitaires sur la chaîne de traitement du signal, d'intégration sur les points d'entrée,
de bout en bout sur le parcours d'analyse — reste à construire. Elle figure parmi les
perspectives.

% =============================================================================
\section{Conclusion}
% =============================================================================

Ce chapitre a décrit la plateforme qui rend les modèles utilisables. L'interface
applicative expose trente points d'entrée documentés automatiquement, protégés par un
double contrôle de rôle et de propriété, et complétés par un cycle de vie des licences
qui permet de révoquer un accès sans détruire de données.

L'archivage sur stockage objet illustre le principe directeur de la réalisation : ne
jamais faire attendre le praticien. Les résultats sont affichés avant la fin du
transfert, et un échec d'archivage ne compromet pas l'exploitation de l'examen.

L'interface utilisateur guide l'analyse en quatre décisions et superpose sur un axe
commun l'hypnogramme, la saturation et les événements respiratoires. Le compte rendu
est composable section par section.

La pile est enfin entièrement conteneurisée et vérifiée par une chaîne d'intégration
continue qui démarre réellement l'application et contrôle le cloisonnement des rôles.
Deux faiblesses ont été signalées sans détour : les routes d'inférence ne sont pas
authentifiées, et la conformité au règlement européen reste incomplète sur le
chiffrement au repos, le droit à l'effacement et les durées de conservation.

La conclusion générale dresse à présent le bilan de l'ensemble du travail.
